% =====================================================================
%   CHAPTER 1: INTRODUCTION
% =====================================================================

\chapter{Introduction}
\label{chap:introduction}

\begin{table}[H]
\caption{History / Revision – Chapter 1}
\label{tab:history-ch1}
\centering
\begin{tabular}{llll}
\toprule
Version & Date & Author & Remarks \\
\midrule
1.0 & January 22, 2026 & Y. Sangale & Initial template draft \\
\midrule
1.0 & April 14, 2026 & Y. Sangale & FPU Specs draft \\
\bottomrule
\end{tabular}
\end{table}

\section{Purpose of Document}
This document defines the first specification of the \gls{riscv} single-instruction-bit microcontroller, covering architectural, electrical, and functional aspects. Revision 1.1 adds the Floating-Point Unit (FPU) specification as an additional design element.

\section{Scope}
The specification outlines the design decisions, \gls{isa} subset, interfaces, and architectural details of the \gls{riscv} \gls{SoC}. This revision extends the scope to include the IEEE 754-2008 compliant double-precision FPU.


\section{Reference Documents}
\begin{itemize}
    \item \gls{riscv} Unprivileged \gls{isa} Specification, Volume I
    \item \gls{riscv} Privileged Architecture Specification
    \item IEEE 754-2008: Standard for Floating-Point Arithmetic
    \item https://docs.riscv.org/reference/isa/unpriv/f-st-ext.html
    \item RWU “Specification Document Preparation Guidelines”
\end{itemize}

